\documentclass[aspectratio=169]{beamer}
\usepackage[utf8]{inputenc}
\usepackage[T1]{fontenc}
\usepackage{amsmath,amssymb}
\usepackage{booktabs}
\usepackage{colortbl}
\usepackage{array}
\usepackage{listings}
\usepackage{xcolor}
\usepackage{lmodern}
\usepackage{graphicx}

%%% ── Colour palette (matching original slides) ────────────────────────────
\definecolor{dmOrange}{RGB}{232,119,34}
\definecolor{dmDarkBlue}{RGB}{20,60,110}
\definecolor{dmGray}{RGB}{55,55,55}
\definecolor{codeBg}{RGB}{233,241,251}
\definecolor{codeKey}{RGB}{0,80,160}
\definecolor{rowHigh}{RGB}{255,243,175}

%%% ── Beamer theme ─────────────────────────────────────────────────────────
\usetheme{default}
\setbeamercolor{frametitle}{fg=dmOrange,bg=white}
\setbeamercolor{title}{fg=dmOrange}
\setbeamercolor{subtitle}{fg=dmDarkBlue}
\setbeamercolor{normal text}{fg=dmGray}
\setbeamercolor{itemize item}{fg=dmOrange}
\setbeamercolor{itemize subitem}{fg=dmOrange}
\setbeamercolor{alerted text}{fg=dmOrange}
\setbeamerfont{frametitle}{size=\Large,series=\bfseries}
\setbeamerfont{title}{size=\LARGE,series=\bfseries}
\setbeamertemplate{navigation symbols}{}
\setbeamertemplate{footline}[frame number]
\setbeamertemplate{itemize item}{\color{dmOrange}\textbullet}
\setbeamertemplate{itemize subitem}{\color{dmOrange}\small$\blacktriangleright$}

%%% ── Listings: Wolfram/Mathematica ────────────────────────────────────────
\lstdefinestyle{wolfram}{
  backgroundcolor=\color{codeBg},
  basicstyle=\ttfamily\footnotesize\color{codeKey},
  keywordstyle=\bfseries,
  commentstyle=\color{gray}\itshape,
  frame=single, framerule=0pt, rulecolor=\color{codeBg},
  xleftmargin=6pt, xrightmargin=6pt,
  breaklines=true, showstringspaces=false,
}

%%% ── Listings: Python / SymPy ─────────────────────────────────────────────
\lstdefinestyle{python}{
  backgroundcolor=\color{codeBg},
  basicstyle=\ttfamily\footnotesize,
  keywordstyle=\color{codeKey}\bfseries,
  commentstyle=\color{gray}\itshape,
  stringstyle=\color{dmOrange},
  frame=single, framerule=0pt, rulecolor=\color{codeBg},
  xleftmargin=6pt, xrightmargin=6pt,
  breaklines=true, showstringspaces=false,
  language=Python,
  morekeywords={symbols,Or,And,Not,Xor,satisfiable,simplify_logic,
                to_dnf,to_cnf,combinations},
}

%%% ── Helper commands ──────────────────────────────────────────────────────
\newcommand{\codelabel}[1]{\textbf{\footnotesize\color{dmDarkBlue}#1}\par\smallskip}
\newcommand{\Iff}{\Leftrightarrow}
\newcommand{\T}{\top}
\newcommand{\F}{\bot}

%%% ── Meta ─────────────────────────────────────────────────────────────────
\title{Equivalence of Propositions}
\subtitle{Introduction to Discrete Mathematics}
\author{}
\date{}

\begin{document}

% ═══════════════════════════════════════════════════════════════════════════
\begin{frame}
  \vspace{1.8cm}
  \titlepage
\end{frame}

% ═══════════════════════════════════════════════════════════════════════════
\begin{frame}{Overview}
  \begin{itemize}
    \item Propositional \textbf{equivalences} are the primary tool of
          simplification in Boolean algebra.
    \item Using equivalences allows the \textbf{analysis of satisfiability}.
    \item Philosophy and reasoning are full of puzzles that can be simplified
          using equivalences.
  \end{itemize}
  \bigskip
  You will learn about
  \textcolor{dmOrange}{\textbf{propositional equivalences}},
  \textcolor{dmOrange}{\textbf{satisfiability}}, and
  \textcolor{dmOrange}{\textbf{normal forms}}.
\end{frame}

% ═══════════════════════════════════════════════════════════════════════════
\section{Satisfiability}

\begin{frame}[fragile]{Satisfiability}
  \begin{itemize}
    \item A \textbf{tautology} is always true:
          $p \lor \lnot p \equiv \top$
  \end{itemize}
  \begin{columns}[T]
    \column{0.47\textwidth}
      \codelabel{Wolfram}
\begin{lstlisting}[style=wolfram]
TautologyQ[p || !p]
(* Out: True *)
\end{lstlisting}
    \column{0.50\textwidth}
      \codelabel{SymPy}
\begin{lstlisting}[style=python]
from sympy import symbols, Or, Not
from sympy.logic.inference import satisfiable
p = symbols('p')
not satisfiable(Not(Or(p, Not(p))))
# True
\end{lstlisting}
  \end{columns}

  \medskip
  \begin{itemize}
    \item A formula is \textbf{satisfiable} if it is true in at least one case.
    \item A \textbf{contradiction} is always false:
          $p \land \lnot p \equiv \bot$
  \end{itemize}
  \begin{columns}[T]
    \column{0.47\textwidth}
      \codelabel{Wolfram}
\begin{lstlisting}[style=wolfram]
SatisfiableQ[p && !p]
(* Out: False *)
\end{lstlisting}
    \column{0.50\textwidth}
      \codelabel{SymPy}
\begin{lstlisting}[style=python]
from sympy import And
bool(satisfiable(And(p, Not(p))))
# False
\end{lstlisting}
  \end{columns}
\end{frame}

% ───────────────────────────────────────────────────────────────────────────
\begin{frame}[fragile]{Satisfiability — Contingency}
  \begin{itemize}
    \item A \textbf{contingency} may be true or false:
          $p \land q$ is satisfiable but not a tautology.
  \end{itemize}
  \begin{columns}[T]
    \column{0.47\textwidth}
      \codelabel{Wolfram}
\begin{lstlisting}[style=wolfram]
SatisfiableQ[p && q]  (* True  *)
TautologyQ[p && q]    (* False *)
\end{lstlisting}
    \column{0.50\textwidth}
      \codelabel{SymPy}
\begin{lstlisting}[style=python]
q = symbols('q')
# Satisfiable? (returns an assignment)
bool(satisfiable(And(p, q)))   # True
# Tautology?
not satisfiable(Not(And(p, q)))# False
\end{lstlisting}
  \end{columns}

  \bigskip
  \begin{block}{Key distinctions}
    \centering
    \begin{tabular}{lccc}
      \toprule
      & \textbf{Satisfiable} & \textbf{Tautology} \\
      \midrule
      Tautology    & \checkmark & \checkmark \\
      Contingency  & \checkmark & $\times$   \\
      Contradiction& $\times$   & $\times$   \\
      \bottomrule
    \end{tabular}
  \end{block}
\end{frame}

% ═══════════════════════════════════════════════════════════════════════════
\section{Main Equivalences}

\begin{frame}{Main Equivalences}
  Two propositions $p$ and $q$ are \textbf{equivalent} iff
  $p \Iff q$ is a tautology.

  \medskip
  \footnotesize
  \renewcommand{\arraystretch}{1.25}
  \begin{tabular}{lll}
    \toprule
    \textbf{Law} & \textbf{Equivalence 1} & \textbf{Equivalence 2} \\
    \midrule
    Identity        & $p \land \T \Iff p$
                    & $p \lor  \F \Iff p$ \\
    Domination      & $p \lor  \T \Iff \T$
                    & $p \land \F \Iff \F$ \\
    Idempotency     & $p \lor  p \Iff p$
                    & $p \land p \Iff p$ \\
    Double Negation & $\lnot(\lnot p) \Iff p$
                    & \\
    Commutativity   & $p \lor  q \Iff q \lor  p$
                    & $p \land q \Iff q \land p$ \\
    Associativity   & $p \lor (q \lor  r) \Iff (p \lor  q) \lor  r$
                    & $p \land(q \land r) \Iff (p \land q) \land r$ \\
    Distributivity  & $p \lor (q \land r) \Iff (p \lor  q)\land(p \lor  r)$
                    & $p \land(q \lor  r) \Iff (p \land q)\lor (p \land r)$ \\
    De Morgan's     & $\lnot(p \land q) \Iff \lnot p \lor  \lnot q$
                    & $\lnot(p \lor  q) \Iff \lnot p \land \lnot q$ \\
    Absorption      & $p \lor (p \land q) \Iff p$
                    & $p \land(p \lor  q) \Iff p$ \\
    Negation        & $p \lor \lnot p \Iff \T$
                    & $p \land \lnot p \Iff \F$ \\
    \bottomrule
  \end{tabular}
  \normalsize

  \medskip
  \textit{Equivalences are tools to simplify compound propositions.}
\end{frame}

% ═══════════════════════════════════════════════════════════════════════════
\section{Simplification Example}

\begin{frame}{Example: Simplifying a Proposition}
  \textbf{Simplify:} ``These potatoes are sweet and salty \textbf{or} they are
  salty, spicy and (salty or spicy).''

  \medskip
  \small
  \begin{align*}
    & (\mathit{sweet} \land \mathit{salty})
      \lor
      \bigl(\mathit{salty} \land \mathit{spicy}
            \land (\mathit{salty} \lor \mathit{spicy})\bigr) \\[4pt]
    \pause
    \Iff\;
    & (\mathit{sweet} \land \mathit{salty})
      \lor (\mathit{salty} \land \mathit{spicy} \land \mathit{salty})
      \lor (\mathit{salty} \land \mathit{spicy} \land \mathit{spicy})
      \quad\text{(Distributivity)} \\[4pt]
    \pause
    \Iff\;
    & (\mathit{sweet} \land \mathit{salty})
      \lor (\mathit{salty} \land \mathit{spicy})
      \lor (\mathit{salty} \land \mathit{spicy})
      \quad\text{(Idempotency)} \\[4pt]
    \pause
    \Iff\;
    & (\mathit{salty} \land \mathit{spicy})
      \lor (\mathit{salty} \land \mathit{sweet})
      \quad\text{(Idempotency)}
  \end{align*}
  \normalsize

  \pause
  \bigskip
  \textbf{Result:} ``These potatoes are salty and spicy, or salty and sweet.''
\end{frame}

% ───────────────────────────────────────────────────────────────────────────
\begin{frame}[fragile]{Example: Automated Simplification}
  \begin{columns}[T]
    \column{0.47\textwidth}
      \codelabel{Wolfram — BooleanMinimize}
\begin{lstlisting}[style=wolfram]
BooleanMinimize[
  (sweet && salty) ||
  (salty && spicy &&
   (salty || spicy))]
(* (salty && spicy) ||
   (salty && sweet) *)
\end{lstlisting}
    \column{0.50\textwidth}
      \codelabel{SymPy — simplify\_logic}
\begin{lstlisting}[style=python]
from sympy import symbols
from sympy.logic.boolalg import (
    And, Or, simplify_logic)

sweet, salty, spicy = symbols(
    'sweet salty spicy')

expr = Or(
    And(sweet, salty),
    And(salty, spicy,
        Or(salty, spicy)))

simplify_logic(expr)
# (salty & spicy) | (salty & sweet)
\end{lstlisting}
  \end{columns}
\end{frame}

% ═══════════════════════════════════════════════════════════════════════════
\section{Conditional Equivalences}

\begin{frame}{Conditional Equivalences}
  All logical operators can be expressed using basic operators
  ($\lnot$, $\land$, $\lor$).

  \medskip
  \renewcommand{\arraystretch}{1.35}
  \begin{tabular}{ll}
    \toprule
    \textbf{Operator} & \textbf{Equivalence} \\
    \midrule
    Conditional
      & $p \to q \;\Iff\; \lnot p \lor q$ \\[4pt]
    Biconditional $\to$ Conditionals
      & $(p \Iff q)
         \;\Iff\;
         (p \to q) \land (\lnot p \to \lnot q)$ \\[4pt]
    Biconditional
      & $(p \Iff q)
         \;\Iff\;
         (p \land q) \lor (\lnot p \land \lnot q)$ \\
    \bottomrule
  \end{tabular}

  \bigskip
  \begin{itemize}
    \item All propositions can be expressed using only $\lnot$, $\land$, $\lor$.
  \end{itemize}
\end{frame}

% ═══════════════════════════════════════════════════════════════════════════
\section{Normal Forms}

\begin{frame}{Disjunctive Normal Form (DNF)}
  A \textbf{DNF} proposition is the disjunction ($\lor$) of one or more
  conjunctions ($\land$).

  \bigskip
  \begin{itemize}
    \item From our example:
          $(\mathit{salty} \land \mathit{spicy})
            \lor (\mathit{salty} \land \mathit{sweet})$
          \quad\textcolor{dmOrange}{is in DNF.}
    \bigskip
    \item Every proposition has a DNF equivalent,
          obtainable by applying equivalence laws.
  \end{itemize}
\end{frame}

% ───────────────────────────────────────────────────────────────────────────
\begin{frame}[fragile]{DNF: BooleanConvert / to\_dnf}
  \begin{columns}[T]
    \column{0.47\textwidth}
      \codelabel{Wolfram}
\begin{lstlisting}[style=wolfram]
BooleanConvert[
  ((a ~Xor~ b ~Xor~ !c ||
    (a||b) && (b||c||d))
    && (!c||b))
  ~Xor~ (!c && a),
  "DNF"]
(* (!a && !b) || (!b && c) *)
\end{lstlisting}
    \column{0.50\textwidth}
      \codelabel{SymPy}
\begin{lstlisting}[style=python]
from sympy import symbols, Xor
from sympy.logic.boolalg import (
    And, Or, Not, to_dnf)

a, b, c, d = symbols('a b c d')

inner = Or(
    Xor(a, b, Not(c)),
    And(Or(a,b), Or(b,c,d)))

expr = Xor(
    And(inner, Or(Not(c), b)),
    And(Not(c), a))

to_dnf(expr, simplify=True)
# (~a & ~b) | (~b & c)
\end{lstlisting}
  \end{columns}
\end{frame}

% ───────────────────────────────────────────────────────────────────────────
\begin{frame}[fragile]{DNF: XOR — ``For me or against me''}
  Translate: ``You are \textbf{either} for me \textbf{or} against me''
  $\;\longrightarrow\; f \oplus a$

  \medskip
  \begin{columns}[T]
    \column{0.47\textwidth}
      \codelabel{Wolfram}
\begin{lstlisting}[style=wolfram]
BooleanConvert[f ~Xor~ a, "DNF"]
(* (a && !f) || (!a && f) *)
\end{lstlisting}
    \column{0.50\textwidth}
      \codelabel{SymPy}
\begin{lstlisting}[style=python]
f, a = symbols('f a')
to_dnf(Xor(f, a), simplify=True)
# (a & ~f) | (~a & f)
\end{lstlisting}
  \end{columns}
\end{frame}

% ───────────────────────────────────────────────────────────────────────────
\begin{frame}[fragile]{DNF: Counting Function}
  \textbf{Problem:} Find the logical expression that is \textbf{true when
  exactly 2 or 3} of $\{a,b,c,d\}$ are true.

  \begin{columns}[T]
    \column{0.47\textwidth}
      \codelabel{Wolfram}
\begin{lstlisting}[style=wolfram]
expr = BooleanCountingFunction[
         {2,3}, {a,b,c,d}];
BooleanConvert@expr  (* full DNF  *)
BooleanMinimize@expr (* minimal   *)
\end{lstlisting}
    \column{0.50\textwidth}
      \codelabel{SymPy}
\begin{lstlisting}[style=python]
from itertools import combinations
from sympy import symbols, And, Or, Not
from sympy.logic.boolalg import (
    to_dnf, simplify_logic)

a, b, c, d = symbols('a b c d')
vs = [a, b, c, d]

def count_fn(ks, vs):
    n = len(vs)
    return Or(*[
        And(*[vs[i] if i in combo
              else Not(vs[i])
              for i in range(n)])
        for k in ks
        for combo in
            combinations(range(n), k)])

expr = count_fn([2, 3], vs)
print(to_dnf(expr, simplify=True))
print(simplify_logic(expr))
\end{lstlisting}
  \end{columns}
\end{frame}

% ───────────────────────────────────────────────────────────────────────────
\begin{frame}[fragile]{Other Normal Forms}
  \begin{columns}[T]
    \column{0.52\textwidth}
      Other normal forms are also possible:

      \medskip
      \footnotesize
      \renewcommand{\arraystretch}{1.3}
      \begin{tabular}{ll}
        \toprule
        \textbf{Form} & \textbf{$x \oplus y$ expressed as \ldots} \\
        \midrule
        DNF  & $(x \land y) \lor (\lnot x \land \lnot y)$ \\
        \rowcolor{rowHigh}
        CNF  & $(\lnot x \lor y) \land (x \lor \lnot y)$ \\
        NAND & $(x \barwedge y) \barwedge (\lnot x \barwedge \lnot y)$ \\
        NOR  & $(\lnot x \downarrow y) \downarrow (x \downarrow \lnot y)$ \\
        \bottomrule
      \end{tabular}
      \normalsize

    \column{0.45\textwidth}
      \codelabel{SymPy — CNF}
\begin{lstlisting}[style=python]
from sympy.logic.boolalg import to_cnf
x, y = symbols('x y')

to_cnf(Xor(x, y), simplify=True)
# (x | y) & (~x | ~y)
\end{lstlisting}
  \end{columns}

  \bigskip
  \begin{itemize}
    \item \textbf{Universal gates}: NAND ($\barwedge$) and NOR ($\downarrow$)
          can each express all binary operations.
  \end{itemize}
\end{frame}

% ═══════════════════════════════════════════════════════════════════════════
\section{CNF to DNF}

\begin{frame}{CNF to DNF Conversion}
  \textbf{Convert} $(\lnot p \lor q) \land (\lnot p \lor \lnot r)
    \land (\lnot q \lor r)$ \textbf{to DNF}:

  \medskip\small
  \begin{align*}
    & (\lnot p \lor q)
      \land (\lnot p \lor \lnot r)
      \land (\lnot q \lor r) \\[4pt]
    \pause
    \Iff\;
    & \bigl(\lnot p \lor (q \land \lnot r)\bigr)
      \land (\lnot q \lor r)
      \quad\text{(Distributivity)} \\[4pt]
    \pause
    \Iff\;
    & \bigl(\lnot p \lor \lnot(\lnot q \lor r)\bigr)
      \land (\lnot q \lor r)
      \quad\text{(De Morgan's)} \\[4pt]
    \pause
    \Iff\;
    & \bigl(\lnot p \land (\lnot q \lor r)\bigr)
      \lor
      \bigl(\lnot(\lnot q \lor r) \land (\lnot q \lor r)\bigr)
      \quad\text{(Distributivity)} \\[4pt]
    \pause
    \Iff\;
    & \lnot p \land (\lnot q \lor r) \;\lor\; \F
      \quad\text{(Negation)} \\[4pt]
    \pause
    \Iff\;
    & \lnot p \land (\lnot q \lor r)
      \quad\text{(Identity)}
  \end{align*}
  \normalsize

  \pause
  \begin{alertblock}{Complexity Note}
    Converting from CNF to DNF is \textbf{NP-hard}
    (a class of computationally difficult problems — see Lesson~14).
  \end{alertblock}
\end{frame}

% ═══════════════════════════════════════════════════════════════════════════
\section{Satisfiability Research}

\begin{frame}{Satisfiability Research}
  \begin{itemize}
    \item Satisfiability remains an \textbf{open problem} at the heart of
          computer science.
    \item Efficiently determining whether \emph{any} logical expression is
          satisfiable can be very difficult.
    \item The propositional satisfiability problem (\textbf{SAT}) is the main
          subject of annual research conferences.
  \end{itemize}

  \bigskip
  \begin{center}
    \large
    \textit{Is there an efficient algorithm for SAT?}\\[4pt]
    \textcolor{dmOrange}{\textbf{P vs.\ NP — one of the Millennium Prize Problems.}}
  \end{center}
\end{frame}

% ═══════════════════════════════════════════════════════════════════════════
\section{Summary}

\begin{frame}{Summary}
  \begin{itemize}
    \item A proposition is either a
          \textcolor{dmOrange}{\textbf{tautology}},
          a \textcolor{dmOrange}{\textbf{contradiction}}, or a
          \textcolor{dmOrange}{\textbf{contingency}}.
    \pause
    \item Many propositional equivalences are used by
          \texttt{simplify\_logic} (SymPy) to simplify Boolean expressions
          automatically.
    \pause
    \item Simplifications are standardised through
          \textcolor{dmOrange}{\textbf{normal forms}}
          (DNF, CNF, NAND, NOR).
    \pause
    \item The previous propositions were \emph{declarative statements}
          — more expressive logics (predicates, quantifiers) come next.
  \end{itemize}
\end{frame}

\end{document}
